\documentclass[10pt,a4paper]{amsart}
\usepackage[margin=19mm]{geometry}
\usepackage{amsmath,amssymb,mathtools}
\usepackage[scr=rsfs,cal=euler]{mathalfa}
\usepackage{xfrac}
\usepackage[unicode,hidelinks,bookmarks=false]{hyperref}
\newcommand{\ie}{i.e.\ }
\newcommand{\cf}{cf.\ }
\newcommand{\resp}{respectively\ }
\newcommand{\Subsection}{\subsection}
\providecommand{\nobreakdash}{\nobreak\hbox{-}}
\setlength{\emergencystretch}{2em}
\multlinegap0pt
\usepackage{mathtools}
\usepackage{amssymb}
\usepackage{tikz}
\newtheorem{remark}{Remark}
\newtheorem{lemma}{Lemma}
\newtheorem{definition}{Definition}
\newcommand{\C}{\mathbb{C}}
\newcommand{\LL}{\mathcal{L}}
\newcommand{\N}{\mathbb{N}}
\renewcommand{\P}{\mathbb{P}}
\newcommand{\Z}{\mathbb{Z}}
\def\bs{\backslash}
\def\ds{\dots}
\newcommand{\g}{\mathfrak{g}}
\newcommand{\id}{\operatorname{id}}
\newcommand{\im}{\operatorname{Im}}
\def\o{\otimes}
\def\om{\omega}
\def\reg{\mathrm{reg}}
\def\si{\sigma}
\def\spn{\mathrm{span}}
\def\ssq{\subseteq}
\def\vac{\mathbf{1}}
\title{Vector Bundles of Coinvariants for Admissible Affine Vertex Operator Algebras}
\author{Victor Alekseev, Jianqi Liu}
\date{Source: arXiv:2609.09154v1 (CC BY 4.0)}
\begin{document}
\maketitle

\noindent\emph{Source and licence.} Vector Bundles of Coinvariants for Admissible Affine Vertex Operator Algebras. Version arXiv:2609.09154v1 (CC BY 4.0), distributed by arXiv under the Creative Commons Attribution 4.0 International licence, http://creativecommons.org/licenses/by/4.0/, as recorded in the arXiv licence metadata for this exact version and shown on the abstract page https://arxiv.org/abs/2609.09154v1. Full licence notice: \texttt{licenses/arxiv-2609.09154-CC-BY-4.0.txt}.\par\smallskip

\noindent\textit{Editorial scope.} Counted unit: the article's lemma lm:DG1 (source0.tex lines 1110--1120). The article's complete original proof of it is delivered with it (source0.tex lines 1122--1214, one \texttt{proof} environment of 93 lines). No mathematics is written, repaired or re-derived: the statement and the proof are the article's own lines, in the article's own notation, and no retained line is substituted or re-flowed.  Retained uncounted as the source-local context the proof needs: lines 595--600; lines 621--626; lines 968--973; lines 980--985; lines 1028--1038; lines 1068--1073; lines 1075--1081; lines 1267--1269.  Nothing was omitted from any retained span, so the package carries no omission record.  No retained line contains a citation, so the wrapper carries no bibliography and no bibliography entry.  No retained line contains a figure environment, an image-inclusion command or drawing code, so no image is delivered and the package declares no asset dependency.  Editorial formatting only, in addition: an AMS \texttt{amsart} preamble that restates the article's own declarations and macros used by the retained lines, namely the environments \texttt{remark}, \texttt{lemma}, \texttt{definition} on separate counters, together with the standard packages those lines need.  The retained environments print in this document as Remark 1, Lemma 1, Definition 1 in order of appearance, whereas the article prints the same items with the numbering of its own full document.  The complete native source of the article as distributed in the arXiv e-print for this exact version is bundled unchanged as \texttt{sources/209784/source0.tex}.
\begin{equation}\label{chiralLieact}
\begin{split}
    \sigma \cdot (w_1 \otimes\cdots\otimes w_n \otimes g) = \sum_{i=1}^n w_1 \otimes\cdots\otimes (\sigma_{P_i} \cdot w_i) \otimes\cdots\otimes w_n \otimes g \\
    = \sum_{i=1}^n \sum_{j} w_1 \otimes\cdots\otimes \left(A_{i,j}\right)^{M^i}_{[j]}\cdot w_i \otimes\cdots\otimes w_n \otimes f_{i,j}\cdot g
\end{split}
\end{equation}

\begin{remark}\label{remark:coordinatefree}
    In particular, the last proposition implies that for any automorphism $\varphi: C\xrightarrow{\simeq} C$ of the smooth curve $C$ such that $\varphi(P_i)=Q_i$ for $1\leq i\leq n$, and for any choice of local coordinates $s_\bullet$ around $Q_\bullet$, there is a vector space isomorphism
\[
[M^\bullet]_{(C,P_\bullet,t_\bullet)}\cong [M^\bullet]_{(C,Q_\bullet,s_\bullet)}.
\]
\end{remark}

\begin{equation}\label{eq:Piaction}
    \begin{aligned}
\si_{P_i}.u&=  \sum_{j,j_1,\ds ,j_n\geq 0} \binom{r}{j}\binom{-m_1}{j_1}\cdots \widehat{\binom{-m_i}{j_i}}\cdots \binom{-m_n}{j_n} \\
 &\cdot P_i^{r-j} (P_i-P_1)^{-m_1-j_1}\cdots \widehat{(P_i-P_{i})}\cdots(P_i-P_n)^{-m_n-j_n} a_{(j+j_1+\ds\hat{j_i}\ds +j_n-m_i)}u.
    \end{aligned}
\end{equation}

\begin{equation}\label{eq:P0action}
    \begin{aligned}
\si_{P_0}.v&=  -\sum_{j_1,\ds ,j_n\geq 0}\binom{-m_1}{j_1}\cdots \binom{-m_n}{j_n} \\
 &\cdot  (-P_1)^{j_1}\cdots (-P_n)^{j_n} a'_{(r-m_1-j_1-\ds -m_n-j_n)}v,
    \end{aligned}
\end{equation}

\begin{lemma}\label{lm:regdegree}
Let $C\bs P_\bullet=\P^1\bs\{P_0,P_1,\ds ,P_n\}$ and let $t_\bullet=(1/z,z-P_1,\ds ,z-P_n)$. Let $M^0,M^1,\ds, M^n$ be admissible $V$-modules attached to $P_0,P_1,\ds,P_n$, with the chiral Lie algebra action described by~\eqref{eq:Piaction} and~\eqref{eq:P0action}. Then for a given section
\[
\si=a\o z^l (dz)^{1-m}\in \mathcal{L}_{\P^1\bs \{P_0,P_1,\ds ,P_n\}}(V)_\reg,\quad 0\leq l\leq 2m-2,
\]
and homogeneous subspaces $M^0(p_0),M^1(p_1),\ds, M^n(p_n)$, with $p_0,p_1,\ds,p_n\geq 0$, of the admissible $V$-modules, we have
\[
\si_{P_0}.M^0(p_0)\ssq M^0(p_0-m+l+1),\quad  \si_{P_i}.M^i(p_i)\ssq \sum_{j\geq 0}M^i(p_i+m-j-1).
\]

\end{lemma}

\begin{definition}\label{def:bottomgen}
An admissible $V$-module $M=\bigoplus_{d\in \N}M(d)$ is \textbf{generated in degree zero} if
\[
M=\spn_\C\left\{a^1_{(m_1)}\cdots a^r_{(m_r)}u\ :\ r\geq 0,\ a^s\in V,\ m_s\in \Z,\ u\in M(0)\right\}.
\]
\end{definition}

\begin{lemma}\label{lm:bottomgen}
Let $V$ be a VOA strongly generated by $V_1$ and let $M$ be an admissible
$V$-module generated in degree zero. Then for every $d\geq 1$,
\begin{equation}\label{eq:bottomgen}
M(d)=\spn_\C\left\{a^1_{(-j_1)}\cdots a^r_{(-j_r)}u\ :\ a^s\in V_1,\ j_s\geq 1,\ \sum\nolimits_{s=1}^r j_s=d,\ u\in M(0)\right\}.
\end{equation}
\end{lemma}

\begin{equation}\label{eq:gsurj}
M^0(0)\o_{U(\g)} M^\bullet (0)\twoheadrightarrow [M^0\o M^\bullet]_{(C,P_\bullet,t_\bullet)},
\end{equation}

\begin{lemma}\label{lm:DG1}
Let $V$ be a VOA of CFT-type strongly generated by $V_1$, let $(C,P_\bullet,t_\bullet)$
be a smooth $(n+1)$-pointed coordinatized curve of genus zero, and let
$M^0,M^1,\ds,M^n$ be admissible $V$-modules, each generated in degree zero
(Definition~\ref{def:bottomgen}). Then the canonical map
\[
M^0(0)\o M^\bullet(0)\longrightarrow [M^0\o M^\bullet]_{(C,P_\bullet,t_\bullet)},\qquad
u^\bullet\mapsto [u^\bullet],
\]
is surjective. No finiteness assumption is made on the spaces $M^i(0)$.
\end{lemma}

\begin{proof}
Write $N:=M^0\o M^\bullet$ and filter it by total degree,
\[
\mathcal F_d N:=\bigoplus_{d_0+\ds+d_n\leq d} M^0(d_0)\o\cdots\o M^n(d_n),\qquad d\in \N,
\]
so that $\mathcal F_0 N=M^0(0)\o M^\bullet(0)$ and $N=\bigcup_{d\geq 0}\mathcal F_d N$. It
suffices to prove
\begin{equation}\label{eq:filtrationstep}
\mathcal F_d N\ssq \mathcal F_{d-1}N+\LL_{C\bs P_\bullet}(V).N \qquad \text{for all } d\geq 1,
\end{equation}
since iterating~\eqref{eq:filtrationstep} gives
$\mathcal F_dN\ssq \mathcal F_0N+\LL_{C\bs P_\bullet}(V).N$ for every $d$, which is the assertion.

Let $d\geq 1$. By Lemma~\ref{lm:bottomgen} the space $\mathcal F_dN$ is spanned by pure
tensors $w^\bullet=w^0\o\cdots\o w^n$ with $w^k\in M^k(d_k)$ homogeneous,
$\sum_k d_k\leq d$, and such that for at least one index $i$ with $d_i\geq 1$ we may write
\[
w^i=a_{(-j)}u^i,\qquad a\in V_1,\ j\geq 1,\ u^i\in M^i(d_i-j)\ \text{homogeneous}.
\]
(If all $d_k=0$ then $w^\bullet\in\mathcal F_0N$ and there is nothing to prove.)

\emph{Construction of a section.} By Remark~\ref{remark:coordinatefree} we may take
\[
(C,P_\bullet,t_\bullet)=\big(\P^1,\ (P_0,P_1,\ds,P_n),\ (1/z,z-P_1,\ds,z-P_n)\big),
\qquad P_0=\infty.
\]
Because $a\in V_1$, a section of the form $\si=a\o f(z)(dz)^{1-1}=a\o f(z)$,
with $f$ a rational function whose poles lie in $P_\bullet$, is an element of
$\LL_{\P^1\bs P_\bullet}(V)$. Set
\[
\si:=a\o\mu,\qquad
\mu:=\begin{cases}
(z-P_i)^{-j}, & 1\leq i\leq n,\\[2pt]
z^{\,j}, & i=0,
\end{cases}
\]
so that in either case $\mu$ has a pole of order exactly $j$ at $P_i$ and no other pole.
This is the only step that uses the genus-zero hypothesis (together with strong generation in degree one):
for $a\in V_1$ the coefficient of $a$ is an ordinary rational \emph{function} rather than a
section of a power of $\om_C$, and on a rational curve a function with one prescribed pole
and no other pole exists for every pole order $j\geq 1$.

\emph{Expansions.} We record the action of $\si$ at each marked point using
\eqref{eq:Piaction} and~\eqref{eq:P0action}, in which $f(z)=z^r\prod_{k=1}^n(z-P_k)^{-m_k}$.
Note first that for $a\in V_1$ formula~\eqref{eq:P0action} simplifies: as in the proof of
Lemma~\ref{lm:regdegree},
\[
-a'_{(l)}=\sum_{s\geq 0}\tfrac{1}{s!}\big(L(1)^sa\big)_{(-l-s)}=a_{(-l)}+\delta_{l,0}\cdot(\text{scalar}),
\]
since $L(1)a\in V_0=\C\vac$, $L(1)^2a=0$ and $\vac_{(-l-1)}=\delta_{l,0}\id$. The
$L(1)$-correction therefore occurs only in the component $l=0$, where it is a scalar, and it
affects none of the degree bounds below. 
%(In the case of interest $L(1)a=0$ in any event;see the discussion preceding~\eqref{eq:gsurj}.)

\emph{Case $1\leq i\leq n$}: here $\mu=(z-P_i)^{-j}$, i.e.\ $r=0$, $m_i=j$ and $m_k=0$ for
$k\neq i$. Since $\binom{0}{s}=\delta_{s,0}$, every binomial factor attached to an index
$\neq i$ collapses, and
\begin{align*}
\si_{P_i}&=a_{(-j)},\\
\si_{P_k}&=\sum_{s\geq 0}\binom{-j}{s}(P_k-P_i)^{-j-s}a_{(s)}\qquad (k\neq i,\ k\geq 1),\\
\si_{P_0}&=\sum_{s\geq 0}\binom{-j}{s}(-P_i)^{s}a_{(j+s)}.
\end{align*}

\emph{Case $i=0$}: here $\mu=z^{j}$, i.e.\ $r=j$ and all $m_k=0$, so
\[
\si_{P_0}=a_{(-j)},\qquad
\si_{P_k}=\sum_{s\geq 0}\binom{j}{s}P_k^{\,j-s}a_{(s)}\quad (1\leq k\leq n).
\]

In both cases $\si$ acts at the distinguished point $P_i$ by the \emph{single} mode
$a_{(-j)}$, raising degree by exactly $j$, and at every other marked point by a combination
of modes $a_{(s)}$ with $s\geq 0$ if that point is finite, resp.\ $a_{(j+s)}$ with $s\geq 0$
if that point is $P_0$. Since $\deg\big(a_{(s)}\big)=-s$ for $a\in V_1$, in all cases
\begin{equation}\label{eq:nonincreasing}
\si_{P_k}.M^k(d_k)\ssq \bigoplus_{e\leq d_k}M^k(e)\qquad (k\neq i);
\end{equation}
that is, $\si$ raises degree at no marked point other than $P_i$.

\emph{Conclusion.} Put $v^\bullet:=w^0\o\cdots\o w^{i-1}\o u^i\o w^{i+1}\o\cdots\o w^n$, an
element of $\mathcal F_{d-j}N$. Because $\LL_{C\bs P_\bullet}(V)$ acts on $N$ one tensor
factor at a time through the expansions $\si_{P_k}$, see~\eqref{chiralLieact}, and because
$\si_{P_i}.u^i=a_{(-j)}u^i=w^i$, we get
\[
\si.v^\bullet-w^\bullet
=\sum_{k\neq i}\left(w^0\o\cdots\o \si_{P_k}.w^k\o\cdots\o u^i\o\cdots\o w^n\right),
\]
and every summand on the right lies in $\mathcal F_{d-j}N\ssq \mathcal F_{d-1}N$ by
\eqref{eq:nonincreasing}, using $j\geq 1$. As
$\si.v^\bullet\in \LL_{C\bs P_\bullet}(V).N$, this proves~\eqref{eq:filtrationstep}.
(Should $\si$ vanish in
$\LL_{C\bs P_\bullet}(V)=H^0(C\bs P_\bullet,\mathscr{V}_C\o\om_C)/\im\nabla$, the same identity
gives $w^\bullet\in\mathcal F_{d-1}N$ directly.)
\end{proof}

\end{document}
